\section{Comparison with adapted FedDualAvg at fixed selected-path uploads}
\label{sec:budget}

We additionally compare coordinate-descent averaging with a deterministic
specialisation of \textsc{FedDualAvg} \citep{yuan2021fco}. This experiment uses
the first $50$ recorded seeds in each simulation design, $5000$--$5049$.
The penalty $\bar\lambda$ is selected as in Section~\ref{sec:sim}, but every
site uses this \emph{common} penalty in both methods. Thus the objective and
its local penalty decomposition are fixed across methods. The
coordinate-descent arm is labelled \textsc{CD-common} to distinguish it
from the locally tuned \textsc{FedAvg} fits in the main study. A converged
pooled Lasso at $\bar\lambda$ supplies an optimisation reference, without
being assigned a federated communication cost.

The adapted \textsc{FedDualAvg} follows Algorithm~3 of \citet{yuan2021fco}
with full participation, sample-size weights, full local empirical
gradients, $h(b)=\|b\|_2^2/2$, $\psi(b)=\bar\lambda\|b\|_1$ and server
learning rate one. Write $f_j(b)=\|y_j-X_jb\|_2^2/(2m_j)$ for the smooth
local loss. At zero-based round $r$ and local step $k$, it forms
$b_{j,r,k}=\soft(z_{j,r,k},\eta(rE+k)\bar\lambda)$ and updates
$z_{j,r,k+1}=z_{j,r,k}-\eta\nabla f_j(b_{j,r,k})$.
The server averages the final dual states with weights $w_j$ and applies
$\soft(\,\cdot\,,\eta(r+1)E\bar\lambda)$ to retrieve its primal iterate. All dual
states start at zero. We report the final server iterate, not the
within-round averaged iterate used in some of that paper's theoretical
bounds. A local step is a full-gradient evaluation from locally cached Gram
and score matrices; it is not a coordinate-descent sweep. These matrices
are not transmitted. We therefore match uploads, not local arithmetic or
wall-clock time.

For each $E\in\{1,5,20\}$, both methods run without early stopping and are
evaluated at $R\in\{5,10,20,50,100\}$ rounds. For adapted
\textsc{FedDualAvg}, the client step is selected separately at each
$(E,R)$ by minimum weighted validation MSE from
$\eta\in\{0.1,0.3,1\}/L_{\max}$, where
$L_{\max}=\max_j\lambda_{\max}(X_j^\top X_j/m_j)$ is calculated from
one scalar supplied by each site. Ties favour the smaller step. Test
responses and true support are not used for tuning. This finite grid and
fixed server step define the implemented comparator; the experiment does
not optimise over all possible \textsc{FedDualAvg} configurations.

One selected-fit round uploads $8Kp=14{,}400$ bytes, so both displayed fits
have the same selected-path vector upload at a given $R$. However, selecting
the \textsc{FedDualAvg} step requires three fits. For a stand-alone budget
$R$, its total upload is $3R(8Kp)+8K(2+3)$ bytes: three vector paths,
the local penalties and Lipschitz constants, and three validation losses
per site. The \textsc{CD-common} total is $R(8Kp)+8K$ bytes. At $R=100$
these are $4218.87$ and $1406.27$ KiB, respectively, versus $1406.25$ KiB
for either selected vector path. Local penalty-grid fits require local
computation but no additional uploads. Downloads, transport overhead and
offline diagnostic evaluation are excluded. The released cost ledger also
accounts for shared setup and reused path prefixes across the whole grid;
stand-alone budget costs must not be added together to represent that
shared run.

We report true-support $F_1$, test MSE, objective difference from the
pooled reference, and the infinity norm of the unit-step proximal-gradient
mapping
$b-\soft(b-\sum_jw_j\nabla f_j(b),\bar\lambda)$.
A zero mapping is equivalent to the Lasso optimality conditions. These
diagnostics distinguish optimisation agreement from support recovery.
Table~\ref{tab:budget} reports the $100$-round comparison and
Figure~\ref{fig:budget} shows the budget trajectories. Full candidate
results and paired Monte Carlo differences are included with the code.

At $100$ rounds, the ranges of mean true-support $F_1$ across the three local-step settings are $0.732$--$0.749$, $0.713$--$0.734$, $0.518$--$0.549$ for adapted \textsc{FedDualAvg} in the independent, correlated and heterogeneous designs, respectively; the corresponding \textsc{CD-common} ranges are $0.256$--$0.298$, $0.235$--$0.282$, $0.234$--$0.285$. These comparisons condition on the selected step-size path; the additional tuning cost stated above remains part of the complete procedure.


The optimisation and selection comparisons point in the same direction:
adapted FedDualAvg has smaller mean objective gaps and higher mean $F_1$
in every displayed design and local-step setting. For the independent
design at $E=1$, its mean gap is $2.7\times10^{-5}$, compared with
$1.695$ for CD-common. This is evidence for using the existing composite
method when accurate minimisation of the common-penalty objective is
required, rather than trying to correct the objective error by changing
CD epochs alone. It is not evidence of uniform superiority across
statistical targets or equal total tuning budgets.

ST and P in Section~\ref{sec:method} are workflow retrofits: ST adds a
small validation exchange, whereas P trains five candidates. This
comparison does not establish their competitiveness with FedDualAvg.
The main study uses local penalties plus server-threshold tuning; this
experiment uses a common penalty and tunes a dual-averaging step.
These different procedures cannot be ranked by combining results across
sections.

\paragraph{Sensitivity to complete tuning cost.}
Supplement S1, Table S2, applies a common total-upload cap of
$1{,}440{,}120$ bytes ($1406.37$~KiB) to the archived runs, counting all
candidate paths and the stand-alone setup and validation scalars above.
For each method and $E$, the largest recorded affordable checkpoint is
selected using cost alone: $R=100$ for CD-common and $R=20$ for adapted
FedDualAvg. Actual uploads are $1406.27$ and $843.87$~KiB, respectively.
The coarse checkpoint grid leaves unused budget, so this is a common-cap
comparison rather than a comparison at equal spending. It was specified
as a sensitivity analysis after the original runs; no new trajectories
or interpolated outcomes are used.
Adapted FedDualAvg retains smaller mean objective gaps and higher mean
$F_1$ in all nine design--$E$ cells; paired mean $F_1$ differences range
from $0.243$ to $0.467$. Thus its advantage over CD-common is not confined
to giving its selected path the same uploads while omitting the other
tuning paths. These results still do not rank ST or P, or match local
computation. Paired differences and their Monte Carlo standard errors
are provided with the code.

\begin{table}[t]
\centering\small
\setlength{\tabcolsep}{3pt}
\caption{Fixed common-penalty comparison after 100 selected-fit upload rounds. Means (Monte Carlo standard errors) over 50 replicates per design. CD is averaging of coordinate-descent iterates with common penalty; DA is adapted FedDualAvg. DA's three-rate search costs 300 vector-upload rounds. $G_\infty$ is the unit-step proximal-gradient mapping residual.}
\label{tab:budget}
\begin{tabular}{llrrrrr}
\toprule
Design & Method & $E$ & Objective gap & $G_\infty$ & $F_1$ & Test MSE \\
\midrule
Independent & CD & 1 & 1.695 (0.051) & 0.213 (0.006) & 0.256 (0.013) & 29.666 (0.564) \\
 & CD & 5 & 1.614 (0.054) & 0.225 (0.006) & 0.295 (0.014) & 30.314 (0.558) \\
 & CD & 20 & 1.642 (0.057) & 0.230 (0.006) & 0.298 (0.014) & 30.380 (0.559) \\
 & DA & 1 & $2.7\times10^{-5}$ ($2.1\times10^{-5}$) & 0.002 (0.001) & 0.732 (0.017) & 28.021 (0.650) \\
 & DA & 5 & $7.9\times10^{-4}$ ($2.1\times10^{-4}$) & 0.013 (0.001) & 0.741 (0.016) & 28.064 (0.650) \\
 & DA & 20 & 0.003 (0.001) & 0.022 (0.002) & 0.749 (0.016) & 28.101 (0.653) \\
\addlinespace
Correlated & CD & 1 & 1.665 (0.037) & 0.206 (0.007) & 0.235 (0.007) & 28.990 (0.518) \\
 & CD & 5 & 1.530 (0.034) & 0.226 (0.008) & 0.278 (0.008) & 29.549 (0.511) \\
 & CD & 20 & 1.548 (0.035) & 0.229 (0.008) & 0.282 (0.007) & 29.588 (0.511) \\
 & DA & 1 & $8.7\times10^{-5}$ ($2.0\times10^{-5}$) & 0.004 (0.000) & 0.713 (0.012) & 27.414 (0.527) \\
 & DA & 5 & 0.001 ($3.7\times10^{-4}$) & 0.014 (0.002) & 0.728 (0.011) & 27.454 (0.529) \\
 & DA & 20 & 0.003 (0.001) & 0.021 (0.003) & 0.734 (0.011) & 27.478 (0.529) \\
\addlinespace
Heterogeneous & CD & 1 & 3.122 (0.086) & 0.395 (0.016) & 0.234 (0.008) & 44.771 (0.822) \\
 & CD & 5 & 2.448 (0.066) & 0.408 (0.016) & 0.277 (0.007) & 44.913 (0.804) \\
 & CD & 20 & 2.475 (0.067) & 0.413 (0.017) & 0.285 (0.007) & 44.997 (0.798) \\
 & DA & 1 & 0.002 (0.001) & 0.016 (0.006) & 0.518 (0.023) & 42.911 (0.937) \\
 & DA & 5 & 0.015 (0.004) & 0.058 (0.006) & 0.539 (0.025) & 42.949 (0.938) \\
 & DA & 20 & 0.039 (0.013) & 0.099 (0.011) & 0.549 (0.028) & 43.024 (0.943) \\
\addlinespace
\bottomrule
\end{tabular}
\end{table}


\begin{figure}[t]
\centering
\includegraphics[width=\textwidth]{figures/budget_comparison.pdf}
\caption{Common-penalty comparison at fixed selected-fit upload budgets.
CD denotes \textsc{CD-common}; DA denotes adapted \textsc{FedDualAvg}.
Error bars are Monte Carlo standard errors over $50$ paired replicates per
design. DA's validation search requires three vector paths; the horizontal
axis excludes those additional tuning paths. Dotted lines show the
same-penalty pooled reference for $F_1$. Local steps differ between methods.}
\label{fig:budget}
\end{figure}
